\PassOptionsToPackage{dvipsnames}{xcolor}
\documentclass[12pt]{beamer}
\usetheme{default}
\usecolortheme{crane}

\usepackage[T1]{fontenc}
\usepackage[slovak]{babel}
\usepackage{amsmath,amssymb}
\usepackage{graphicx,array,booktabs,alltt,listings,tikz}
\usetikzlibrary{arrows.meta,positioning}
\setbeamertemplate{footline}[frame number]
\beamertemplatenavigationsymbolsempty
\setbeameroption{hide notes}

\def\blue#1{\textcolor{Cerulean}{#1}}
\def\green#1{\textcolor{LimeGreen}{#1}}
\newcommand\qsp{\hskip.25em plus.08333em minus.083333em}
\newcommand\dash{\unskip\qsp\hbox{--\kern-.25em--}\qsp}
\providecommand{\uv}[1]{\quotedblbase #1\textquotedblleft}
\newcommand{\dfn}[1]{\emph{#1}}
\newcommand{\quant}[2]{\mathop{#1#2\/}\nolimits}
\newcommand{\A}{\quant\forall}
\newcommand{\E}{\quant\exists}
\newcommand{\proves}{\vdash}
\newcommand{\nproves}{\nvdash}
\newcommand{\entails}{\vDash}
\newcommand{\nentails}{\nvDash}
\newcommand{\limpl}{\rightarrow}
\newcommand{\lequiv}{\leftrightarrow}
\newcommand{\Equiv}{\mathrel{\Leftrightarrow}}
\newcommand{\emptyclause}{\Box}
\newcommand{\RULE}[2]{\ensuremath{\begin{array}{c}#1\\\hline#2\end{array}}}
\newcommand{\bigO}{\mathcal{O}}
\newcommand{\Lang}[1][]{\mathcal{L}_{\text{#1}}}

% Based on the Lean listings in Logika 2, p1/p1.tex.
\lstdefinelanguage{Lean}{
  keywords={variable,theorem,def,by,intro,obtain,exact,Prop,example,import,
    List,Nat,Type,fun,let,have,show,if,then,else,sorry},
  comment=[l]{--},
  literate={∃}{$\exists$}{1} {∀}{$\forall$}{1}
    {∧}{$\land$}{1} {∨}{$\lor$}{1} {¬}{$\neg$}{1}
    {→}{$\to$}{1} {↔}{$\leftrightarrow$}{1}
    {⟨}{$\langle$}{1} {⟩}{$\rangle$}{1}
    {ℕ}{$\mathbb{N}$}{1} {∈}{$\in$}{1}
    {≤}{$\le$}{1} {≥}{$\ge$}{1} {≠}{$\ne$}{1}
    {×}{$\times$}{1} {₁}{$_1$}{1} {₂}{$_2$}{1},
  basicstyle=\ttfamily\small,
  keywordstyle=\color{blue}\bfseries,
  commentstyle=\color{gray}\itshape,
  showstringspaces=false,
  columns=fullflexible,
  breaklines=true,
  keepspaces=true
}

\title{Logika, vyjadrovacia sila\\a formálna verifikácia}
\author{Ján Mazák}
\institute{FMFI UK Bratislava}
\date{}

\begin{document}
\frame{\titlepage}

\section{Formálne teórie a hranice rozhodovania}
% Existing frames are copied from the sources noted at the start of each frame.

% Adapted from Logika 1, pr12.pdf p. 4 (informational final lecture).
\begin{frame}
  Cieľom tejto prednášky je umožniť vám získať predstavu o matematickej logike
  a niektorých jej významných výsledkoch.
  Prednáška zároveň naznačí, prečo dotazovacie jazyky majú obmedzenú vyjadrovaciu silu.
  K tvrdeniam uvedeným v prezentácii pristupujte opatrne,
  podané sú v zjednodušenej forme a občas sú zamlčané technickejšie predpoklady.

  \bigskip

  Prednáška je informačná; obsah nebude na skúške.
\end{frame}

% Terminology from Logika 1, pr07.pdf p.17, 7-kvantifikatory/2-syn.tex;
% add function symbols as used in pr12, frame Peanova aritmetika.
\begin{frame}{Jazyk logiky 1.~rádu}
  \begin{itemize}
    \item \alert{Logické symboly}: spojky $\lnot,\land,\lor,\limpl$,
      kvantifikátory $\forall,\exists$ a~rovnosť~$=$.
    \item Indivíduové premenné $x,y,\ldots$ a~pomocné symboly (zátvorky, čiarka).
    \item \alert{Mimologické symboly} si volíme podľa toho, o~čom chceme hovoriť:
      \begin{itemize}
        \item predikátové symboly pre vzťahy, napr. binárny~$R$;
        \item funkčné symboly pre funkcie, napr. unárny~$f$;
        \item indivíduové konštanty pre konkrétne objekty, napr.~$c$.
      \end{itemize}
  \end{itemize}
  Počet argumentov symbolu je jeho \emph{arita}.

  \medskip
  Interpretácia určí konkrétne objekty, vzťahy a~funkcie.
  Napr. nad $\mathbb N$ môžeme $c$ interpretovať ako~$0$,
  $f$ ako $n\mapsto n+1$ a~$R$ ako~$\le$.
\end{frame}

% New mathematical examples requested for theory, model and entailment.
% Partial orders are already covered in the compulsory discrete mathematics.
\begin{frame}{Teória: čiastočné usporiadanie}
  \alert{Teória} je množina uzavretých formúl (bez voľných premenných) vo zvolenom jazyku.
  Napr. jazyk má jeden binárny predikátový symbol~$R$ a~teória~$T$ obsahuje formuly
  \begin{align*}
    &\forall x\ R(x,x)
      &&\text{reflexivita},\\
    &\forall x\forall y\ (R(x,y)\land R(y,x)\limpl x=y)
      &&\text{antisymetria},\\
    &\forall x\forall y\forall z\
      (R(x,y)\land R(y,z)\limpl R(x,z))
      &&\text{tranzitivita}.
  \end{align*}
  To, či je zápis formulou, je čisto syntaktická otázka.
  Zatiaľ neurčujeme množinu prvkov ani konkrétny význam~$R$.
  O~pravdivosti formuly hovoríme až po voľbe domény
  a~interpretácie mimologických symbolov.
\end{frame}

% Logika 1, pr07.pdf p.29, 7-kvantifikatory/3-sem.tex;
% retain the definition and add the interpretation of function symbols.
\begin{frame}{Štruktúra}
  Nech $\Lang$ je jazyk logiky prvého rádu.
  \alert{Štruktúrou} pre jazyk~$\Lang$
  nazývame dvojicu $\mathcal M=(D,i)$, kde
  \begin{itemize}
    \item $D$ je ľubovoľná \alert{neprázdna} množina
      nazývaná \emph{doména} štruktúry~$\mathcal M$;
    \item $i$ je zobrazenie, nazývané \emph{interpretačná funkcia}, ktoré
      \begin{itemize}
        \item každej indivíduovej konštante~$c$ jazyka~$\Lang$
          priraďuje prvok $i(c)\in D$;
        \item každému predikátovému symbolu~$P$ s~aritou~$n$
          priraďuje reláciu $i(P)\subseteq D^n$;
        \item každému funkčnému symbolu~$f$ s~aritou~$n$
          priraďuje funkciu $i(f)\colon D^n\to D$.
      \end{itemize}
  \end{itemize}
  Rovnosť sa interpretuje ako rovnosť prvkov domény~$D$.
\end{frame}

\begin{frame}{Model teórie}
  \alert{Model teórie~$T$} je štruktúra, v~ktorej sú splnené všetky formuly z~$T$.
  Píšeme $\mathcal M\models T$.
  \begin{itemize}
    \item $D=\mathbb N$, $i(R)={\le}$: model našej teórie~$T$.
    \item $D=\mathcal P(\{1,2\})$, $i(R)={\subseteq}$: tiež model~$T$. ($\mathcal P(A)$ označuje množinu všetkých podmnožín~$A$.)
    \item $D=\mathbb Z$, $i(R)={<}$: nie je model~$T$, neplatí reflexivita.
  \end{itemize}

  \medskip
  Rovnaká teória môže mať rôzne modely, konečné aj nekonečné.
\end{frame}

% Axiom terminology is introduced only after theories and models.
% One illustrative formula per theory, not a list of their axioms.
\begin{frame}{Teórie pre konkrétne oblasti matematiky}
  Pri opise konkrétnej oblasti matematiky si vyberáme konkrétnu teóriu
  a~jej formulám hovoríme \alert{axiómy}.

  \medskip
  \alert{Peanova aritmetika (PA)}: aritmetika prirodzených čísel.
  Napr. jedna z~axióm:
  \[
    \forall x\ (x+0=x).
  \]

  \alert{ZFC}: teória množín. Napr. axióma extenzionality:
  \[
    \forall A\forall B\ \bigl(\forall x\,(x\in A\leftrightarrow x\in B)
      \limpl A=B\bigr).
  \]
  Množiny s~rovnakými prvkami sa rovnajú.
  \note{V ZFC premenné A, B aj x označujú objekty domény, teda množiny;
    ide o kvantifikátory prvého rádu. Uvedené sú iba dva príklady formúl,
    nie celé axiomatizácie PA či ZFC. Vzorec PA je zo zdrojových slidov
    Logika 1, pr12, rámec Peanova aritmetika.}
\end{frame}

\begin{frame}{Pravdivosť v modeli a vyplývanie z teórie}
  Pre uzavretú formulu~$F$ zápis
  \alert{$\mathcal M\models F$} znamená, že $F$ je splnená v~štruktúre~$\mathcal M$ (\uv{pravdivá pre konkrétnu interpretáciu mimologických symbolov}).

  \medskip
  \alert{$T\models F$}: formula~$F$ platí v~\alert{každom modeli} teórie~$T$.

  \medskip
  Pre našu teóriu čiastočného usporiadania uvažujme
  \[
    F=\forall x\forall y\ (R(x,y)\lor R(y,x)).
  \]
  \vspace*{-5mm}
  \begin{itemize}
    \item Pre $D=\mathbb N$ a~$i(R)={\le}$ platí~$F$.
    \item Pre $D=\mathcal P(\{1,2\})$ a~$i(R)={\subseteq}$ neplatí~$F$:
      množiny $\{1\}$ a~$\{2\}$ sú neporovnateľné.
  \end{itemize}
  Preto $T\not\models F$ a~$T\not\models\lnot F$ --- formula $F$ je od teórie $T$ nezávislá.
\end{frame}

% A finite-sequence definition, relative to chosen inference rules.
% The example of an inference rule is copied from Logika 2, p1/p1.tex.
\begin{frame}{Formálny dôkaz a dokázateľnosť}
  Zvoľme odvodzovacie pravidlá dokazovacieho kalkulu.
  \alert{Dôkaz formuly~$F$ z~teórie~$T$} je konečná postupnosť
  $F_1,\ldots,F_n=F$, kde každá formula~$F_i$
  \begin{itemize}
    \item patrí do~$T$, alebo
    \item vznikne z~predchádzajúcich formúl použitím odvodzovacieho pravidla.
  \end{itemize}
  Napr. pravidlo modus ponens:
  \[
    \frac{A\qquad A\limpl B}{B}.
  \]
  $T\vdash F$: existuje takýto formálny dôkaz~$F$ z~$T$.
  \note{Pravidlá môžu mať aj nulový počet premís. Modus ponens je iba
    príklad, nie celý kalkul. Iné kalkuly môžu dôkaz zapisovať stromom
    či termom; podstatné sú presné lokálne podmienky jednotlivých krokov.}
\end{frame}

% Adapted from Matematická logika, p1.pdf p. 10; original p1/p1.tex.
\begin{frame}{Rezolvencia}
\textbf{Rezolvencia} operuje nad klauzulami (disjunkciami literálov)
a odvodzuje nové klauzuly, kým nenájde spor (prázdnu klauzulu).

\[
\frac{(A \lor C_1\lor C_2) \qquad (\lnot A \lor B_1\lor B_2)}
     {C_1\lor C_2 \lor B_1\lor B_2}
\]

\textbf{Dôkaz sporom:} pri dôkaze~$F$ z~$T$ pridáme $\lnot F$
a~po prepise do klauzúl zachovávajúcom splniteľnosť
odvodíme prázdnu klauzulu.

\textbf{Korektnosť:} rezolventa vyplýva z rodičovských klauzúl.

\medskip
Používame tu malú sadu pravidiel na odvodenie prázdnej klauzuly,
ale dôkaz zostáva postupnosťou algoritmicky ľahko overiteľných krokov,
rovnako ako v~mnohých ďalších dokazovacích systémoch,
ktoré logici počas dejín vytvorili.

\end{frame}

% Adapted from Matematická logika, p1.pdf p. 11; original p1/p1.tex.
\begin{frame}{Rezolvencia: príklad dôkazu}
Dokážeme nesplniteľnosť množiny
\[
  \{ A\lor B,\qquad \lnot A,\qquad \lnot B\}.
\]
\begin{enumerate}
  \item Rezolvujeme $A\lor B$ s $\lnot A$, dostaneme $B$.
  \item Rezolvujeme $B$ s $\lnot B$, dostávame prázdnu klauzulu,
        tá je nesplniteľná.
\end{enumerate}
\end{frame}

% Adapted from Logika 1, pr10.pdf p. 37; original 10-def-korekt/7-unifikacia.tex.
\begin{frame}{Rezolvencia v logike 1.~rádu: unifikácia}
Čo možno usúdiť z nasledujúcich dvoch formúl?
\[
  \forall y\,P(f(y),y),
  \qquad
  \forall x\,(\lnot P(x,d)\lor R(x))
\]
Dosadenie popisujeme pomocou substitúcie:
\[
  \sigma=\{x\mapsto f(d),\ y\mapsto d\}.
\]
Po substitúcii~$\sigma$ majú komplementárne literály rovnaké argumenty
predikátu (preto $\sigma$ nazývame \emph{unifikátor}):
$P(f(d),d)$ a $\lnot P(f(d),d)$.

Jedným z dôsledkov uvedených dvoch formúl je teda $R(f(d))$.
\end{frame}

% Adapted from Logika 1, pr12.pdf p. 6; original 12-advanced/1-druhoradova.tex.
\begin{frame}{Pravdivosť a dokázateľnosť}
  Uvažujme jazyk logiky 1.~rádu a v ňom teóriu $T$ a formulu $F$.
  \begin{block}{Veta o~korektnosti logiky 1.~rádu}
    Ak $T\proves F$, tak $T\entails F$.
  \end{block}
  \begin{block}{Veta o~úplnosti logiky 1.~rádu}
    Ak $T\entails F$, tak $T\proves F$.
  \end{block}

  Úplnosť nie je samozrejmá: $T\models F$ hovorí o~všetkých modeloch,
  aj o~tých s~nespočítateľnou doménou, ktorej prvky nemožno
  algoritmicky enumerovať.

  Pre platnú formulu musí existovať dôkaz. Pre neplatnú musí existovať protipríklad (model, v ktorom neplatí).
  Ani jedno z toho však nemusíme vedieť algoritmicky nájsť.
\end{frame}

% New short orientation slide, based on Logika 1, pr12.pdf p. 21.
\begin{frame}{Nerozhodnuteľné problémy}
\begin{itemize}
  \item \alert{Rozhodnuteľný problém:} existuje algoritmus, ktorý vždy skončí
        a odpovie správne áno alebo nie.
  \item \alert{Čiastočne rozhodnuteľný problém:} algoritmus vždy potvrdí
        odpoveď áno, ak je správna; pri odpovedi nie môže bežať navždy.
  \item \alert{Ani čiastočne rozhodnuteľný:} neexistuje ani takýto
        jednostranne zaručený algoritmus.
\end{itemize}

\medskip
Príklady nerozhodnuteľných problémov:
\begin{itemize}
  \item \emph{Problém zastavenia programu:} zastaví sa daný program pre daný vstup?
  \item \emph{Ekvivalencia dvoch programov v~C:} majú pre každý vstup rovnaký
    výsledok?
\end{itemize}
\end{frame}

% Theorem statements adapted from Logika 1, pr12.pdf pp. 16 and 21;
% original 12-advanced/2-nerozhodnutelnost.tex. Proofs omitted.
\begin{frame}{Platnosť formuly vo FOL}
\begin{block}{Veta}
  Platnosť formuly v prvorádovej logike je nerozhodnuteľná.
\end{block}

Z~korektnosti a~úplnosti: formula je platná práve vtedy, keď má dôkaz.
Systematicky enumerujeme všetky dôkazy; ak je formula platná,
jej dôkaz skôr či neskôr nájdeme.
Tak dostávame \alert{čiastočný rozhodovací postup}.

\medskip
Na hľadaní dôkazov sú založené automatické dokazovače.
Rezolvencia s~unifikáciou sa vo FOL v~praxi osvedčila ako účinná metóda automatického dokazovania.
Opakovane odvodzujú pomocou heuristík nové a nové dôsledky (v súlade s pravidlami) a dúfajú, že raz trafia kompletný dôkaz.

\medskip
Otázka, či veta FOL \emph{nie je} platná, nie je ani čiastočne rozhodnuteľná.
\note{Uvažujeme efektívne zadaný jazyk. Úplné systematické hľadanie
  zaručuje nájdenie každého existujúceho dôkazu; táto teoretická záruka
  sa nesmie zamieňať so zárukou praktického heuristického či časovo
  obmedzeného behu dokazovača.}
\end{frame}

% Restored from Logika 1, pr12.pdf p.21, 2-nerozhodnutelnost.tex.
% Satisfiability statement from p.21; monadic FOL from p.22 as a brief curiosity.
\begin{frame}{Splniteľnosť v prvorádovej logike}
  \begin{block}{Veta}
    Splniteľnosť v~prvorádovej logike je nerozhodnuteľná;
    nie je ani čiastočne rozhodnuteľná.
  \end{block}

  \medskip
  \emph{Kuriozita:} monadická prvorádová logika (MFOL) nemá funkčné symboly
  (ani konštanty) a~jej predikáty majú len jeden argument.

  Platnosť formuly v~MFOL je rozhodnuteľná.
  Ak pridáme čo len jeden predikát arity~2, platnosť formuly už
  nie je rozhodnuteľná.
\end{frame}


% New early connection: domain independence and safe relational queries.
% Zdroj: Abiteboul, Hull, Vianu, Foundations of Databases, 5.4 a 6.3.
\begin{frame}{Bezpečné databázové dotazy}
  Dotaz je \alert{doménovo nezávislý}, ak rozšírenie domény
  o~ďalšie hodnoty bez zmeny tabuliek nezmení jeho výsledok.

  \begin{itemize}
    \item $\lnot\mathrm{Zamestnanec}(x)$ nie je doménovo nezávislý:
      vracia aj hodnoty mimo databázy. Pri rozšírení domény pribudnú
      ďalšie výsledky.
    \item $\mathrm{Osoba}(x)\land\lnot\mathrm{Zamestnanec}(x)$
      vyberá iba osoby uložené v~tabuľke, ktoré nie sú zamestnancami.
      Rozšírenie domény výsledok nezmení.
  \end{itemize}

  Relačná algebra a~relačné jadro SQL umožňujú iba
  \emph{bezpečné dotazy}: výsledok vieme vypočítať z~konečných dát,
  bez prehľadávania ľubovoľnej domény.

  \alert{Pre FOL nemožno algoritmicky rozhodnúť, či je zadaný dotaz
  doménovo nezávislý}, preto bezpečnosť zaručujeme obmedzením
  povolených zápisov (priamo syntaxou dotazovacieho jazyka).
  \note{Uvažujeme klasické relačné dotazy v množinovom modeli,
    bez všeobecných funkcií a rekurzie; príklady majú unárne tabuľky.
    Doména obsahuje aj konštanty uvedené v dotaze. Doménová nezávislosť
    je sémantická vlastnosť; bezpečné syntaktické obmedzenia ju zaručujú.
    Nerozhodnuteľnosť splniteľnosti FOL neznamená, že nevieme
    vyhodnotiť formulu na konkrétnej konečnej doméne.}
\end{frame}

% Shortened from Logika 1, pr12.pdf p. 30; 4-neuplnost.tex.
% Pri enumerovaných axiómach možno ich použitie overiť podľa priloženého
% počtu krokov enumerácie; samotné členstvo v množine axióm nemusí byť rozhodnuteľné.
\begin{frame}{Efektívna axiomatizovateľnosť}
  Teória je \alert{efektívne axiomatizovateľná}, ak spĺňa tieto
  požiadavky na algoritmickú prácu:
  \begin{itemize}
    \item \emph{Symboly:} vieme algoritmicky rozpoznať ich zápis.
    \item \emph{Termy a~formuly:} vieme rozhodnúť, či je zápis správny.
    \item \emph{Axiómy:} algoritmus ich postupne všetky vypíše;
      môže ich byť nekonečne veľa.
    \item \emph{Odvodzovacie pravidlá:} vieme algoritmicky overiť
      ich použitie.
  \end{itemize}
  Potom môžeme systematicky enumerovať všetky dôkazy
  a~dokázateľné formuly.
\end{frame}

% Adapted from Logika 1, pr12.pdf p. 31; original 12-advanced/4-neuplnost.tex.
\begin{frame}{Sémantická a~negačná úplnosť}
  Doterajší pojem úplnosti, tzv. \alert{sémantická úplnosť} formálneho systému,
  hovoril o~tom, že pravdivé veci možno dokázať:\\
  Pre každú teóriu $T$ a~uzavretú formulu $F$, ak $T\entails F$, tak $T\vdash F$.

  Zaujímavá je však aj \alert{negačná úplnosť} teórie $T$:
  pre každú uzavretú formulu $F$ (v~jazyku tejto teórie)
  platí $T\vdash F$ alebo $T\vdash \lnot F$.
\end{frame}

% Theorem statement adapted from Logika 1, pr12.pdf p. 34;
% original 12-advanced/4-neuplnost.tex. Proof omitted.
\begin{frame}{Prvá G\"odelova veta o~neúplnosti}
  \begin{block}{Veta}
    Prvorádová teória~$T$ nemôže mať súčasne všetky štyri vlastnosti:
    \begin{enumerate}
      \item je konzistentná (nedokazuje spor);
      \item je efektívne axiomatizovateľná;
      \item obsahuje dostatočnú aritmetiku (napr. PA);
      \item je negačne úplná: pre každú uzavretú formulu~$F$ jej jazyka
        platí $T\vdash F$ alebo $T\vdash\lnot F$.
    \end{enumerate}
  \end{block}
  V princípe: zmysluplné teórie nikdy nehovoria o~všetkom, vždy zostanú
  nejaké tvrdenia, ktoré z~teórie nemožno ani dokázať, ani vyvrátiť.
  Matematiku ako celok nemožno plne axiomatizovať.
\end{frame}

% Theorem statement adapted from Logika 1, pr12.pdf p. 38;
% original 12-advanced/4-neuplnost.tex. Proof omitted.
\begin{frame}{Druhá G\"odelova veta o~neúplnosti}
  \begin{block}{Veta}
    Ak $S$ je bezosporný efektívne axiomatizovateľný
    formálny systém obsahujúci aritmetiku (napr. PA),
    jeho bezospornosť (pri obvyklej formalizácii)
    nemožno dokázať v~rámci~$S$.
  \end{block}
  Čiže dôkaz treba spraviť neformálnou metaúvahou alebo v~rámci iného \uv{silnejšieho} formálneho systému.
  Ako však dokážeme bezospornosť tohto druhého systému? \dots
  \note{Štandardná formalizácia vlastných dôkazov a bezospornosti;
    dostatočná aritmetika pre obvyklé podmienky odvoditeľnosti.}
\end{frame}

\section{Vyjadrovacia sila a rekurzia}
% New orientation after Godel's theorems: expressive limits of FOL.
\begin{frame}{Hranice vyjadrovacej sily FOL}
  \begin{itemize}
    \item Vo FOL kvantifikujeme cez prvky domény;
      predikáty (v~databáze tabuľky) sú pevne dané.
    \item Chceli by sme sa opýtať: \uv{Existuje v~databáze iná
      tabuľka disjunktná s~tabuľkou~$R$?}
      Pri pevnej schéme to vo FOL vieme rozpísať cez jednotlivé tabuľky.
      Chceli by sme však vyberať tabuľku premennou.
    \item Ak pridáme premenné pre predikáty a~kvantifikujeme ich,
      dostaneme \alert{logiku 2.~rádu}. Pri plnej sémantike premenné
      prechádzajú všetky možné relácie, nielen existujúce tabuľky.
    \item Ani dosiahnuteľnosť v~grafe sa nedá vyjadriť prvorádovými formulami, ako si čochvíľa dokážeme.

  \end{itemize}
\end{frame}

% Logika 1, pr12.pdf p.7; 1-druhoradova.tex.
% Preserve the question; compactness and its use in the proof follow.
\begin{frame}{Ako vyjadriť tranzitívny uzáver}
  Majme binárnu reláciu reprezentovanú predikátom~$R$.
  Pre každé $n\ge 1$ označme $R_n(a,b)$ formulu
  \[
    \E{x_1}\E{x_2}\cdots\E{x_n}
    \bigl(R(a,x_1)\land R(x_1,x_2)\land\dots\land R(x_n,b)\bigr).
  \]
  Ďalej nech $R_0=R$ a~nech
  \[
    R^*(a,b)\quad\Equiv\quad
      \hbox{existuje~$n$ také, že platí $R_n(a,b)$}.
  \]
  Naše $R^*$ vyjadruje tranzitívny uzáver relácie~$R$.

  \bigskip
  Dá sa však $R^*$ vyjadriť prvorádovou formulou?
  \note{Ide o dosiahnuteľnosť cestou kladnej dĺžky. Rekurzívny SQL dotaz
    už bol v p5a; neopakovať ho. Kompaktnosť sama osebe nedokazuje
    nevyjadriteľnosť pri obmedzení na konečné štruktúry.}
\end{frame}

% Compactness from Logika 1, pr12.pdf p.6, 1-druhoradova.tex.
% Added equivalent formulation and a short finite/infinite-model discussion.
\begin{frame}{Kompaktnosť logiky 1.~rádu}
  \begin{block}{Veta o~kompaktnosti}
    Každá nesplniteľná teória má konečnú podmnožinu, ktorá je nesplniteľná.
  \end{block}
  Ekvivalentne: ak má každá konečná podmnožina teórie~$T$ model,
  potom má model aj celá teória~$T$.
  \begin{itemize}
    \item Teória~$T$ pritom môže byť nekonečná.
    \item Výsledný model nemusí byť konečný. Napr. požiadavky \uv{aspoň $n$ prvkov} pre každé $n\ge1$:
  každý konečný výber má konečný model, všetky spolu iba nekonečný.
  \end{itemize}
\end{frame}

% Shortened from Logika 1, pr12.pdf p.8, 1-druhoradova.tex.
% One-slide proof over arbitrary structures, with the finite-model caveat visible.
% Zdroj pre konečný prípad: Libkin, Elements of Finite Model Theory, kapitoly 3 a 4.
\begin{frame}{Tranzitívny uzáver: dôkaz kompaktnosťou}
  Predpokladajme, že $R^*$ sa dá vyjadriť prvorádovou formulou
  vo všetkých štruktúrach. Uvažujme teóriu
  \[
    T=\{R^*(a,b)\}\cup\{\lnot R_n(a,b):n\ge0\}.
  \]
  \begin{enumerate}
    \item $T$ nemá model: ak $R^*(a,b)$, tak pre nejaké~$n$ platí $R_n(a,b)$.
    \item Každá konečná podmnožina~$T$ má model:
      pre dostatočne veľké~$n$ vyhovuje cesta z~$a$ do~$b$ dĺžky~$n+2$.
  \end{enumerate}
  Tieto dve pozorovania sú v~spore s~kompaktnosťou FOL.

  \medskip
  \alert{Formulou logiky 1.~rádu tranzitívny uzáver nevieme vyjadriť.}

  \medskip
  (Nevyjadriteľnosť sa dá dokázať aj pre konečné modely.)
  \note{Kompaktnosť dáva model, nie nevyhnutne konečný model. Ak by
    kandidátna formula mala správne vyjadrovať dosiahnuteľnosť iba na
    konečných grafoch, nekonečný model nemusí dať spor. Konečný prípad
    vyžaduje iný argument (napr. Ehrenfeuchtove-Fraisseho hry).
    Tento konečný výsledok potom ukazuje, prečo na všeobecnú
    dosiahnuteľnosť nestačí prvorádové relačné jadro SQL a potrebujeme
    rozšírenie, napr. rekurziu/fixpoint. Rekurzívny dotaz neopakovať.}
\end{frame}

% Adapted from Logika 2, p1.pdf pp.37-38, frame Databázy.
% Connection to the recursive reachability query already covered in p5a.
\begin{frame}{Databázy: prečo potrebujeme rekurziu}
  \begin{itemize}
    \item Relačné jadro SQL bez rekurzie zodpovedá FOL
      (množinový model, bezpečné dotazy).
    \item Na dosiahnuteľnosť v grafe teda potrebujeme rozšírenie jazyka:
      rekurzia počítajúca \alert{najmenší pevný bod}
      (\emph{least fixed point}, LFP) prekonáva túto hranicu FOL.
    \item Najmenší pevný bod vypočítame už známou naivnou
      alebo seminaivnou evaluáciou.
  \end{itemize}
  \note{Nejde o celé SQL s agregáciami, NULL a všeobecnými funkciami.
    Opisujeme lineárnu rekurziu pre dosiahnuteľnosť: k už nájdenej ceste
    pripojíme jednu hranu a odstránime duplicity. Obe evaluácie počítajú
    tú istú najmenšiu reláciu obsahujúcu hrany a uzavretú na toto pravidlo.
    Seminaivná evaluácia vynechá rozvíjanie starých dvojíc.
    Aj na konečných grafoch dosiahnuteľnosť nie je vyjadriteľná vo FOL;
    tento fakt sme uviedli bez dôkazu. LFP je jedno užitočné rozšírenie;
    netvrdíme, že každý rekurzívny SQL dotaz je monotónny alebo skončí.}
\end{frame}

% Aggregation extends the relational core; ignore NULL and empty-group conventions.
% MIN/MAX assume an available order; COUNT parity uses ordinary arithmetic.
% SUM can aggregate the constant 1 per row, so it subsumes COUNT.
% Source: Libkin, Expressive power of SQL,
% https://homepages.inf.ed.ac.uk/libkin/papers/tcs-icdt01.pdf
\begin{frame}{Databázy: agregácie a~FOL}
  \begin{itemize}
    \item \texttt{MIN}, \texttt{MAX}: pri dostupnom usporiadaní
      nepridávajú nič navyše oproti FOL.
    \item \texttt{SUM}: pridáme do logiky číselný term
      $\sum_{x:\varphi(x)}t(x)$: súčet hodnôt $t(x)$ pre všetky
      riadky spĺňajúce $\varphi(x)$.
      Počet sčítancov závisí od dát, zložitá sémantika.
    \item \texttt{COUNT}: osobitný prípad súčtu, kde $t(x)=1$;
      s~aritmetikou vyjadrí aj paritu, ktorú obyčajná FOL nevyjadrí.
    \item \texttt{AVG}: nič navyše oproti \texttt{SUM + COUNT}.
  \end{itemize}

  Ani s~agregáciami nemožno vyjadriť tranzitívny uzáver.
  % Libkin, Expressive Power of SQL, Theorem 3 and the discussion in Section 8.
  % Arithmetic is on the separate numeric sort; graph vertices are unordered.
  % Do not generalize this bound to arbitrary order/arithmetic on vertex IDs.
\end{frame}

\section{Databázové jazyky a optimalizácia}
% New slide, agreed language-design tradeoff.
% Zdroje: https://www.postgresql.org/docs/current/xfunc-optimization.html
% https://www.postgresql.org/docs/current/queries-with.html
% https://www.postgresql.org/docs/current/plpython-trigger.html
\begin{frame}{Prečo nepoužiť C alebo Python nad tabuľkami?}
  \begin{itemize}
    \item Bezpečné relačné jadro bez rekurzie a~všeobecných funkcií
      sa nad konečnou databázou vždy vyhodnotí.
      Pre ľubovoľný program nevieme všeobecne rozhodnúť ani to, či skončí.
    \item C aj Python sa v~DBMS používať dajú.
      Funkcia v \uv{neobmedzenom} jazyku použitá ako súčasť dotazu je však pre optimalizátor len čierna skrinka.
    \item Vyjadrovaciu silu pridávame tam, kde je užitočná či nutná:
      napr. rekurzia pre hierarchické dáta.
    \item Zložitejšie kontroly konzistencie, ktoré sa ťažko vyjadrujú
      v~SQL, možno napr. v~PostgreSQL napísať ako trigger functions
      v~Pythone.
  \end{itemize}
  \note{Python API zostavujúce relačný plán môže zachovať výhody
    deklaratívneho dotazovania. Rozhoduje sprístupnená štruktúra, nie názov
    jazyka. Moderné SQL s rozšíreniami nemusí skončiť. Neprisudzovať
    relačnej algebre rozhodnuteľnosť všetkých otázok o dotazoch.}
\end{frame}

% New slide: finite equivalence and the semidecidability reversal.
% Zdroje: https://doi.org/10.1007/978-1-4615-9385-0_13
% https://pmc.ncbi.nlm.nih.gov/articles/PMC7324044/
\begin{frame}{Sú dva stromy relačnej algebry ekvivalentné?}
  Optimalizácia je náročná už pre relačnú algebru:
  môžeme jeden strom nahradiť iným bez zmeny výsledku?
  \begin{itemize}
    \item Ekvivalentné stromy vrátia rovnaký výsledok na
      \emph{každej konečnej databáze}.
    \item Pre plnú relačnú algebru (vrátane rozdielu) je ekvivalencia
      \alert{nerozhodnuteľná}; nie je ani čiastočne rozhodnuteľná.
      % Trakhtenbrotova veta.
  \end{itemize}
  Optimalizátor používa dokázané pravidlá prepisovania;
  nemusí rozpoznať všetky ekvivalentné dvojice stromov.
  \note{Klasický množinový model bez NULL a dodatočných integritných
    obmedzení; rovnaká výstupná schéma. Plná algebra zahŕňa rozdiel.
    Pri FOL nad všetkými štruktúrami je čiastočne rozhodnuteľná
    ekvivalencia; nad konečnými databázami je čiastočne rozhodnuteľná
    neekvivalencia. Tento rozdiel možno spomenúť ústne.
    Pre konjunktívne dotazy s rovnosťami v množinovom modeli
    je ekvivalencia rozhodnuteľná.}
\end{frame}

% New slide. Illustration of interacting choices, not a reduction.
\begin{frame}{Optimalizácia pri pevnom strome}
  Jednoduchšia úloha: zafixujeme strom v relačnej algebre, všetky joiny rátame ako \alert{merge join},
  vyberáme si iba poradie atribútov, podľa ktorého sa relácie pri výpočte joinu usporiadajú (zhoda sa posudzuje na celej $n$-tici spoločných atribútov):
  \begin{center}
  \begin{tikzpicture}[x=1cm,y=1cm,
    every node/.style={font=\small},
    edge/.style={draw,->}]
    \node (top) at (0,1.7) {join na $b$};
    \node (left) at (-2.5,.85) {join na $(a,b)$};
    \node (right) at (2.5,.85) {join na $(b,c)$};
    \node (r) at (-3.5,0) {$R$};
    \node (s) at (-1.5,0) {$S$};
    \node (t) at (1.5,0) {$T$};
    \node (u) at (3.5,0) {$U$};
    \path[edge] (r)--(left);
    \path[edge] (s)--(left);
    \path[edge] (t)--(right);
    \path[edge] (u)--(right);
    \path[edge] (left)--(top);
    \path[edge] (right)--(top);
  \end{tikzpicture}
  \end{center}
  \begin{itemize}
    \item Naľavo $(b,a)$ zaručí zoradenie podľa~$b$ aj pre horný join;
      $(a,b)$ túto záruku nedáva.
    \item Napravo analogicky: $(b,c)$, nie $(c,b)$.
      Existujúci index môže zvýhodniť poradie nevhodné pre horný join.
  \end{itemize}
  Aj výber zoradení pre pevný strom merge joinov je vo všeobecnosti
  \alert{NP-ťažký}. Úloha je však rozhodnuteľná.

  %  {\scriptsize \href{https://arxiv.org/abs/cs/0611094}{Guravannavar a~kol. (2006)}, oddiel~4.1.}
  % \note{Strom určuje závislosti, nie časový rozvrh nezávislých vetiev.
  %   Veta 4.1 dokazuje NP-ťažkosť zjednodušeného modelu: vyberáme permutácie
  %   stĺpcov v rastúcom strome a maximalizujeme súčet dĺžok spoločných
  %   prefixov na hranách. Nejde o tvrdenie, že tento malý príklad je ťažký,
  %   ani že každé prideľovanie operátorov pevnému stromu je NP-ťažké.
  %   Konečné efektívne ohodnotiteľné množiny kandidátov možno prehľadať:
  %   NP-ťažké nie je nerozhodnuteľné. Menšie prípady možno riešiť presne;
  %   všeobecnú polynomiálnu metódu by vylúčilo P != NP.
  %   Odhadovaná cena nie je zaručený skutočný čas.}
\end{frame}

% Copied from Logika 2, p1.pdf p.32, frame Logika 2. radu.
% Omit reading assignment; qualify full second-order semantics.
\begin{frame}{Logika 2. rádu}
  Syntax a~sémantiku logiky 1.~rádu možno rozšíriť o premenné pre predikáty, ktoré je dovolené kvantifikovať.
  Dostaneme tak logiku 2.~rádu (SO).\\[3mm]

  Ak sa pri kvantifikovaní predikátov obmedzíme na predikáty arity 1, dostaneme monadickú logiku 2. rádu (MSO).
  Takéto predikáty de facto zodpovedajú množinám prvkov domény, čiže v MSO možno vyjadriť bežné veci z teórie grafov apod.\\[3mm]

  Pri plnej sémantike:
  \alert{Použitie formúl 2.~rádu zvyšuje vyjadrovaciu silu jazyka, ale oproti FOL stratíme kompaktnosť a~úplnosť}.
\end{frame}

% Adapted from Logika 1, pr12.pdf p.24, 3-descriptive.tex.
% Fagin: https://kenclarkson.org/RonFaginPaper/genspec.pdf
% Immerman-Vardi: https://people.cs.umass.edu/~immerman/descriptiveComplexity.html
\begin{frame}{Deskriptívna zložitosť}
  Pri klasickom pohľade na zložitosť algoritmov skúmame
  počet krokov v~závislosti od veľkosti vstupu.

  Alternatívny pohľad: akú zložitú logickú formulu potrebujeme
  na popis daného problému?

  Príklady rôzne zložitých formúl:
  \begin{itemize}
    \item existential SOL:\quad
      $\E{X_1}\cdots\E{X_n}\ F$, kde $F$ je formula FOL
    \item universal SOL:\quad
      $\A{X_1}\cdots\A{X_n}\ F$, kde $F$ je formula FOL
  \end{itemize}
  \alert{Formula je pevná, vstupom sú dáta (napr. graf).}
\end{frame}

% Copied from Logika 1, pr12.pdf p.25, 3-descriptive.tex.
% Add the fixed-formula qualification to the original complexity bound.
\begin{frame}{Deskriptívna zložitosť}
  Existencia trojuholníka v~grafe sa dá vyjadriť vo FOL:
  \begin{align*}
    &\E x\E y\E z\ \big(V(x)\land V(y)\land V(z)\land {}\\
    &\qquad{}\land x\not\doteq y\land y\not\doteq z\land z\not\doteq x
      \land E(x,y)\land E(y,z)\land E(z,x)\big)
  \end{align*}

  Predikát $V$~popisuje vrcholy, $E$~hrany.
  Formulu vyhodnocujeme na konkrétnom grafe,
  neoverujeme jej všeobecnú platnosť.

  Pre $k$~kvantifikátorov vo formule preveríme $O(|V|^k)$ možností.
  \alert{Formula je pevná, rastie graf.}
  Preto tento problém patrí do triedy~P.
\end{frame}

% Copied from Logika 1, pr12.pdf p.26, 3-descriptive.tex.
% Schematic formula: the ellipsis omits the other colours and the partition condition.
\begin{frame}{Deskriptívna zložitosť}
  Existencia 3-farbenia grafu: NP-úplný;\\
  nevyjadriteľné vo FOL, vyjadriteľné v existential SOL:
  \[
    \E R\E G\E B
    \bigl(\A x\A y\ (E(x,y)\limpl\lnot R(x)\lor\lnot R(y))\land\dotsb\bigr)
  \]
  ($R$, $G$, $B$ sú predikáty vyjadrujúce jednotlivé farby).\\

  Za doménu zoberieme vrcholy grafu.
  Množiny farieb nedeterministicky uhádneme;
  vyhodnotenie pevnej prvorádovej časti už trvá polynomiálny čas.
  Preto tento problém patrí do triedy NP.
\end{frame}

% Copied from Logika 1, pr12.pdf p.27, 3-descriptive.tex.
% Built-in order/arithmetic is needed for the uniform AC^0 characterization.
% Do not assert PH is strictly contained in PSPACE: that separation is open.
% The SO fixed-point operator is of higher order than the FO fixed-point operator.
% Sources: https://people.cs.umass.edu/~immerman/descriptiveComplexity.html
% https://cgi.cse.unsw.edu.au/~eptcs/paper.cgi?FICS2024.7.pdf=
\begin{frame}{Deskriptívna zložitosť}
  Konečné usporiadané vstupy; formula je pevná, rastú dáta.
  \begin{itemize}
    \item FOL s~aritmetikou na pozíciách: trieda zložitosti AC$^0$
      (vlastná podmnožina LOGSPACE)
    \item FOL + tranzitívny uzáver: non-deterministic LOGSPACE
    \item FOL + least fixed point operator: P
  \end{itemize}
  \begin{itemize}
    \item existential SOL: NP
    \item universal SOL: co-NP
    \item SOL: PH (obsahuje NP aj coNP)
    \item SOL + least fixed point operator: EXPTIME
  \end{itemize}
  V~poslednom riadku ide o~pevný bod vyššieho rádu.
\end{frame}

\section{Lean a správnosť špecifikácie}
% Source: Logika 1, pr12.pdf pp. 48-51, 55-59.
% Reused from 12-advanced/6-lean.tex and lean_code/; #print and
% standalone specification slide omitted by agreement.

\begin{frame}{Lean: jazyk a interaktívny dokazovací systém}
Lean je dôkazový asistent založený na teórii typov.
\begin{itemize}
    \item Na dokázanie formuly treba skonštruovať term zodpovedajúceho typu.
    \item Lean vie verifikovať existujúce dôkazy (termy).
    \item Lean tiež ponúka tzv. \uv{taktiky}, ktoré pomáhajú automatizovať konštrukciu dôkazu
      (napr. mu možno dať pokyn \uv{sprav rozbor prípadov}, \uv{nejako skús použiť túto lemu},
      \uv{aplikuj pravidlá násobenia v poli}).
    \item \texttt{mathlib} je rastúca knižnica formálne verifikovaných matematických tvrdení.
\end{itemize}
\end{frame}

\begin{frame}[fragile]{Lean: dôkaz priamou konštrukciou termu}
$A,B,C$ sú tvrdenia. Dôkaz tvrdenia~$A$ je v~Leane hodnota typu~$A$.
\begin{lstlisting}[language=Lean,basicstyle=\footnotesize\ttfamily]
-- Direct proof term
theorem my_term (A B C : Prop) :
    (A ∧ C) → (B → A) :=
  fun h_ac : A ∧ C =>
    fun h_b : B =>
      h_ac.left
\end{lstlisting}
Funkcia dostane dôkaz $A\land C$, potom dôkaz~$B$
a~vráti dôkaz~$A$. Tým dokazuje celú implikáciu.
\end{frame}

\begin{frame}[fragile]{Lean: dôkaz pomocou taktík}
Ten istý dôkazový term môžeme zostrojiť pomocou taktík:
\begin{lstlisting}[language=Lean,basicstyle=\footnotesize\ttfamily]
theorem my_tactic_proof (A B C : Prop) :
    (A ∧ C) → (B → A) := by
  intro h_ac
  obtain ⟨h_a, h_c⟩ := h_ac
  intro h_b
  exact h_a
\end{lstlisting}
{\small Výroky ako typy, dôkazy ako termy: Curryho-Howardova korešpondencia.}
\end{frame}



% Source: Logika 1, pr12.pdf p. 55; copied with the original sorry holes.
\begin{frame}[fragile]{Lean: kód s~dôkazom správnosti}
Lean je aj funkcionálny programovací jazyk.

\medskip
\textbf{Úloha:} napísať funkciu pre maximum z~neprázdneho zoznamu
a~dokázať, že výsledok patrí do zoznamu a~je najväčší.

\medskip
\textbf{Možnosť 1:} návratový typ obsahuje výsledok aj dôkaz správnosti.
\begin{lstlisting}[language=Lean,basicstyle=\scriptsize\ttfamily]
def find_max : (L : List ℕ) → L ≠ [] →
        { x // x ∈ L ∧ ∀ y ∈ L, y ≤ x }
  | [], h => by contradiction
  | [a], _ => ⟨a, by sorry⟩
  | a :: b :: tail, _ =>
      let ⟨m, ⟨hm_mem, hm_max⟩⟩ := find_max (b :: tail) (by simp)

      if h : m ≤ a then ⟨a, by sorry⟩
      else ⟨m, by sorry⟩
\end{lstlisting}
{\small Dôkazové časti sú vynechané (\texttt{sorry}).}
\end{frame}

\begin{frame}[fragile]{Lean: samostatný dôkaz správnosti}
\textbf{Možnosť 2:} funkcia vracia bežnú hodnotu a~jej správnosť
vyjadríme samostatnou vetou.

Veta požaduje, aby výsledok patril do zoznamu a~bol najväčší.
\begin{lstlisting}[language=Lean,basicstyle=\scriptsize\ttfamily]
def find_max : (L : List ℕ) → L ≠ [] → ℕ := sorry

theorem find_max_correct : (L : List ℕ) → (h : L ≠ []) →
  find_max L h ∈ L ∧ ∀ y ∈ L, y ≤ find_max L h := by sorry
\end{lstlisting}
{\small Náčrt deklarácií; implementácia a~dôkaz sú vynechané.}
\end{frame}

% Logika 1, pr12.pdf pp.58-59: original theorem declarations.
\begin{frame}[fragile]{Lean: dokázateľne správny kód}
  Čo presne znamená \alert{rozdeliť zoznam na párne a nepárne prvky}?
  Zachytávajú tieto podmienky všetko, čo od programu chceme?
  {\small Ukážka špecifikácie; dôkazy sú vynechané.}
\begin{lstlisting}[language=Lean,basicstyle=\scriptsize\ttfamily]
theorem split_parity_length (xs : List Nat) :
  (split_parity xs).1.length + (split_parity xs).2.length = xs.length := by sorry

theorem split_parity_all_even (xs : List Nat) :
  ∀ x ∈ (split_parity xs).1, Even x := by sorry

theorem split_parity_all_odd (xs : List Nat) :
  ∀ x ∈ (split_parity xs).2, Odd x := by sorry

theorem split_parity_mem (xs : List Nat) (x : Nat) :
  x ∈ xs ↔ x ∈ (split_parity xs).1 ∨ x ∈ (split_parity xs).2 := by sorry
\end{lstlisting}
\end{frame}

\begin{frame}{Dôkaz je len taký dobrý ako špecifikácia}
Chýba zachovanie násobností. Požiadavky na dĺžku, paritu a~členstvo dovoľujú
\[
  [2,4,4] \longmapsto ([2,2,4],[]).
\]
\begin{itemize}
    \item Súčet dĺžok výstupných zoznamov je tri.
    \item Prvý výstup obsahuje iba párne čísla a druhý iba nepárne.
    \item Vstup a výstup obsahujú tie isté rôzne hodnoty, ale nie tie isté násobnosti.
\end{itemize}
\end{frame}

\begin{frame}[fragile]{Dôkaz je len taký dobrý ako špecifikácia}
Chýba stabilita: aj keď zachováme násobnosti, slabá špecifikácia povoľuje meniť poradie.

Popri podmienkach pre paritu a~dĺžku požadujeme:
\begin{lstlisting}[language=Lean,basicstyle=\scriptsize\ttfamily]
theorem split_parity_stable (xs : List Nat) :
  (split_parity xs).1.Sublist xs ∧
  (split_parity xs).2.Sublist xs := by sorry
\end{lstlisting}
Podzoznam vznikne vynechaním prvkov bez zmeny poradia.
Spolu s~paritou a~súčtom dĺžok tým zaručíme aj správne násobnosti.

\medskip
Stabilita je veľmi dôležitá aj pri testovaní: pri nedeterministickom
výstupe je porovnanie s~jediným očakávaným zoznamom nespoľahlivé.

V~praxi môže premiešanie časovo zoradených záznamov pokaziť
napr. výpočet zmien medzi susednými meraniami.
\end{frame}



\section{Solvery}
% From logic algorithms to practical, independently checkable computation.
% Davis, Logemann, Loveland (1962); Chaff (2001):
% https://www.princeton.edu/~chaff/publication/DAC2001v56.pdf
% Conflict-driven clause learning: https://www.princeton.edu/~chaff/publication/iccad2001_final.pdf
% Zelkova uses a solver portfolio including Z3; the daily invocation count
% applies to Zelkova, not to Z3 alone and not only to S3.
% https://assets.amazon.science/56/3c/7d75bca84c56b9b1ed1999beec0c/semantic-based-automated-reasoning-for-aws-access-policies-using-smt.pdf
% https://www.microsoft.com/en-us/research/project/z3-3/
% Certificate verification in Lean:
% https://lean-lang.org/doc/api/Lean/Elab/Tactic/BVDecide.html
\begin{frame}{Logické solvery v~praxi}
  \begin{itemize}
    \item Logické solvery sa stali jednými z~najrýchlejších nástrojov
      na mnohé praktické aj matematické problémy.
    \item Základ algoritmov pre SAT poznáme od~60.~rokov;
      praktický prelom prišiel okolo~2000. CDCL (Conflict-Driven Clause Learning) sa pri backtrackingu
      učí z~konfliktov nové klauzuly a~v~prehľadávanom priestore odrezáva celé podstromy.
    \item Amazon používa Z3 od Microsoftu pri kontrole prístupu
      k~S3 (cloudové úložisko). Ide o~SMT solver, ktorého princíp si o~chvíľu vysvetlíme.
    \item Výsledkom môže byť overiteľný certifikát: ohodnotenie (SAT)
      alebo dôkaz nesplniteľnosti (UNSAT).
    \item Certifikáty využijeme v~Leane pri overovaní kódu a~výpočtov,
      v~priemysle aj matematickom výskume.
  \end{itemize}
\end{frame}

% Reused and lightly edited from Logika 1, pr06.pdf.

% Original: Logika 1, pr06.pdf, p. 6, "Problém SAT".
\begin{frame}{Problém SAT}
  \begin{itemize}
    \item \alert{Má boolovská formula ohodnotenie, v ktorom je splnená?}
    \item Pre danú boolovskú formulu hľadáme priradenie hodnôt
      pravda/nepravda jej premenným, pri ktorom je formula pravdivá.
      Existuje také priradenie?
    \item Vstupom býva formula v~CNF: konjunkcia klauzúl,
      ktorú už poznáme z~rezolvencie.
    \item SAT solver nájde vyhovujúce priradenie alebo rozhodne,
      že žiadne neexistuje.
    \item \emph{Predspracovanie:} rezolvenčná eliminácia znižuje
      počet premenných aj za cenu väčších klauzúl,
      čo môže urýchliť následné hľadanie.
    \item Hoci je SAT NP-úplný problém, solvery sú prekvapujúco rýchle, neraz aj pre tisíce či milióny premenných.
  \end{itemize}
\end{frame}

% Original: Logika 1, pr06.pdf p.8, plus graph applications from p.95.
% Practical strength across domains is the main point, not solver internals.
\begin{frame}{Praktické využitie SAT solverov}
  \begin{itemize}
    \item Verifikácia hardvéru: existuje vstup, na ktorom sa obvod
      pred optimalizáciou a~po nej správa odlišne?
    \item Softvérové balíky: existuje kompatibilný výber verzií?
    \item Existuje farbenie grafu najviac $k$~farbami?
    \item Existuje v~grafe hamiltonovská kružnica?
  \end{itemize}
  Každý problém najprv preložíme na boolovskú formulu.
  Do SAT možno zakódovať prakticky ľubovoľnú konečnú kombinatorickú úlohu;
  veľkosť prekladu však nemusí byť polynomiálna
  (pri bežných kombinatorických úlohách však je).
\end{frame}



% Original: Logika 1, pr06.pdf, p. 109, "Čo je SMT solver?".
\begin{frame}{Čo je SMT solver?}
  \alert{SMT} (Satisfiability Modulo Theories) kombinuje logické spojky
  s~podmienkami zo zvolených teórií, napr. aritmetiky.

  Pre $x,y\in\mathbb Z$ hľadáme riešenie:
  \[
    (x=0\lor y=0)\land x+y=5\land x>y.
  \]

  \begin{itemize}
    \item Podporované teórie: lineárna aritmetika nad celými
      a~reálnymi číslami, funkcie bez určenej interpretácie,
      polia, bitové vektory.
    \item Úplné rozhodovanie sa týka vybraných teórií a~fragmentov,
      nie ľubovoľnej matematiky.
  \end{itemize}
\end{frame}



% SAT/theory cooperation: de Moura and Bjorner, Tractability and Modern
% Satisfiability Modulo Theories Solvers, section 1.1.
% https://www.microsoft.com/en-us/research/wp-content/uploads/2016/02/nbjorner-tractability.pdf
\begin{frame}{Ako pracuje SMT solver?}
  \begin{itemize}
    \item Základom je SAT solver, ale atómami môžu byť aj matematické
      tvrdenia, napr. $x+y=5$ alebo $x>y$.
      SAT časť ich význam nepozná; berie ich ako boolovské premenné.
    \item SAT solver navrhuje, ktoré atómy majú byť pravdivé
      a~ktoré nepravdivé, aby vyhovel logickým spojkám.
    \item Špecializované moduly pre jednotlivé teórie priebežne kontrolujú,
      či sú tieto voľby zlučiteľné s~doterajšími podmienkami
      v~príslušnej teórii.
    \item Napr. modul pre lineárnu aritmetiku odhalí spor
      v nerovnostiach a~oznámi ho SAT solveru.
      Ten sa z~konfliktu poučí a~hľadá iné ohodnotenie.
  \end{itemize}
\end{frame}


\section{Verná formalizácia}
\begin{frame}{Dokázali sme zamýšľané tvrdenie?}
Pri počítačovom dokazovaní aj databázových dotazoch je kľúčová
\alert{formalizácia zadania z~prirodzeného jazyka}.
Môže ju vytvoriť človek alebo LLM.

\begin{itemize}
  \item Algoritmické strojové overenie či výpočet
    \alert{nezaručí, že zápis vystihuje pôvodné zadanie.}
  \item \uv{Nech $n$ je prvý krok, v~ktorom sa program $P$ zastaví.}
    Tento zápis predpokladá, že taký krok existuje.
  \item Overiť tento skrytý predpoklad znamená vyriešiť
    problém zastavenia.
\end{itemize}

\alert{Už overenie vernosti formalizácie môže byť nerozhodnuteľné,
dokonca ťažšie než problém zastavenia.}
% TODO make more precise
V~princípe k~ľubovoľne ťažkému výpočtovému problému tak existuje aspoň tak ťažký formalizačný problém.
Ani Lean, ani LLM nás od toho nezachránia.

\vfill
{\tiny Bastounis, Circelli, Hansen (2026),
\href{https://arxiv.org/abs/2610.08144}{Navier-Stokes lost in translation}.}
\note{Príklad sa týka skrytých existenčných predpokladov a vernosti prekladu.
  Samotný zápis symbolickej či nekonštruktívnej definície nevyžaduje výpočet
  jej hodnoty. Pridanie predpokladu alebo predvolenej hodnoty môže zmeniť
  pôvodné zadanie. Článok ukazuje neobmedzenú výšku v aritmetickej/SCI
  hierarchii; túto hierarchiu tu nedefinujeme.}
\end{frame}


\end{document}
